% Calcul mental 6e : compléments à 100 puis à 360 (préparation des angles).
% Diaporama à défilement automatique : 25 s par question
% (5 overlays de 5 s, \transduration -> entrée /Dur de chaque page du PDF).
% Titre, « Fin » et correction : pas de durée (le professeur avance à la main).
% Avancement automatique respecté par Okular, Adobe Reader (plein écran)...
% mais pas par les visionneuses des navigateurs.
\documentclass[aspectratio=169,12pt]{beamer}
\usepackage[T1]{fontenc}
\usepackage[utf8]{inputenc}
% babel french non installé sur cette machine : compilé sans babel
\usepackage{lmodern}
\usepackage{tikz}
\usepackage{array}

\setbeamertemplate{navigation symbols}{}
\setbeamercolor{background canvas}{bg=white}
\setbeamercolor{frametitle}{fg=black,bg=white}
\setbeamerfont{frametitle}{size=\LARGE,series=\bfseries}

\definecolor{barre}{RGB}{40,110,200}
\definecolor{urgent}{RGB}{200,40,40}

% Barre de compte à rebours : 5 étapes de 5 s (25 s, 20 s, ..., 5 s)
\newcommand{\compteur}{%
  \foreach \k in {1,...,5}{%
    \only<\k>{%
      \pgfmathsetmacro{\reste}{6-\k}%
      \pgfmathtruncatemacro{\sec}{5*(6-\k)}%
      \begin{tikzpicture}
        \draw[gray!60,line width=0.8pt,rounded corners=3pt] (0,0) rectangle (10,0.6);
        \ifnum\k<5
          \fill[barre,rounded corners=3pt] (0,0) rectangle ({10*\reste/5},0.6);
        \else
          \fill[urgent,rounded corners=3pt] (0,0) rectangle ({10*\reste/5},0.6);
        \fi
        \node[anchor=west,font=\Large\bfseries] at (10.3,0.3) {\sec\,s};
      \end{tikzpicture}%
    }%
  }%
}

% \question{numéro}{total visé}{nombre donné}
\newcommand{\question}[3]{%
  \begin{frame}
    \transduration<1-5>{5}% 5 overlays x 5 s = 25 s
    \frametitle{Complète pour faire #2}
    \begin{center}
      {\Large Question #1 / 20}\par\vspace{1.2cm}
      {\fontsize{42}{50}\selectfont\bfseries\mbox{#3 + \dots{} = #2}}\par\vspace{1.5cm}
      \compteur
    \end{center}
  \end{frame}
}

\begin{document}

% --- Titre : pas de durée, le professeur lance avec une touche ---
\begin{frame}
  \begin{center}
    {\fontsize{40}{48}\selectfont\bfseries Calcul mental}\par\vspace{0.8cm}
    {\LARGE 20 questions — 25 s par question}\par\vspace{1cm}
    {\Large Compléments à 100, puis à 360\par\vspace{0.2cm}
      \small(360° = un tour complet)}\par\vspace{1cm}
    {\Large\bfseries Écris seulement le résultat,\\[0.2cm] numérote de 1 à 20.}
  \end{center}
\end{frame}

% --- Série 1 : compléments à 100 ---
\question{1}{100}{70}
\question{2}{100}{40}
\question{3}{100}{35}
\question{4}{100}{85}
\question{5}{100}{64}
\question{6}{100}{37}
\question{7}{100}{8}
\question{8}{100}{91}
\question{9}{100}{43}
\question{10}{100}{19}

% --- Série 2 : compléments à 360 ---
\question{11}{360}{300}
\question{12}{360}{180}
\question{13}{360}{270}
\question{14}{360}{90}
\question{15}{360}{200}
\question{16}{360}{330}
\question{17}{360}{245}
\question{18}{360}{172}
\question{19}{360}{318}
\question{20}{360}{87}

% --- Fin : pas de durée ---
\begin{frame}
  \begin{center}
    {\fontsize{50}{60}\selectfont\bfseries Fin}\par\vspace{1cm}
    {\LARGE Posez les stylos !}
  \end{center}
\end{frame}

% --- Correction : pas de durée ---
\newcommand{\rep}[1]{\textcolor{urgent}{\textbf{#1}}}

\begin{frame}
  \frametitle{Correction — compléments à 100}
  \begin{center}\LARGE
    \renewcommand{\arraystretch}{1.25}
    \begin{tabular}{>{\bfseries}r l @{\hspace{0.6cm}} >{\bfseries}r l}
      1. & 70 + \rep{30} = 100 &  6. & 37 + \rep{63} = 100 \\
      2. & 40 + \rep{60} = 100 &  7. & 8 + \rep{92} = 100 \\
      3. & 35 + \rep{65} = 100 &  8. & 91 + \rep{9} = 100 \\
      4. & 85 + \rep{15} = 100 &  9. & 43 + \rep{57} = 100 \\
      5. & 64 + \rep{36} = 100 & 10. & 19 + \rep{81} = 100 \\
    \end{tabular}
  \end{center}
\end{frame}

\begin{frame}
  \frametitle{Correction — compléments à 360}
  \begin{center}\Large
    \renewcommand{\arraystretch}{1.25}
    \begin{tabular}{>{\bfseries}r l @{\hspace{0.6cm}} >{\bfseries}r l}
      11. & 300 + \rep{60} = 360  & 16. & 330 + \rep{30} = 360 \\
      12. & 180 + \rep{180} = 360 & 17. & 245 + \rep{115} = 360 \\
      13. & 270 + \rep{90} = 360  & 18. & 172 + \rep{188} = 360 \\
      14. & 90 + \rep{270} = 360  & 19. & 318 + \rep{42} = 360 \\
      15. & 200 + \rep{160} = 360 & 20. & 87 + \rep{273} = 360 \\
    \end{tabular}
  \end{center}
\end{frame}

\end{document}
